\section{Study Population}

The study population comprises adults with histologically confirmed, resectable or
borderline-resectable, KRAS G12-mutated pancreatic ductal adenocarcinoma (PDAC)
who are eligible for daraxonrasib under broad pan-KRAS criteria and whose anatomy
is compatible with the on-premises eight-arm robotic system for those who may
receive the experimental arm. This is a multicenter, randomized, controlled,
open-label study conducted across eight academic high-volume hepatopancreatobiliary
(HPB) centers. Target enrollment is approximately 245 participants screened to
yield 220 randomized 1:1 (110 per arm), where the modest screen-to-randomization
margin allows for up to 10 percent non-evaluability while preserving the planned
statistical power (\S9). Participants are allocated to Arm A (experimental:
perioperative daraxonrasib at the recommended Phase 2 dose plus the LLM-directed
robotic Whipple) or Arm B (active control: modified FOLFIRINOX \cite{Conroy2018}
plus institutional-standard high-volume pancreaticoduodenectomy). Figure 10
presents the randomized two-arm CONSORT participant flow from screening through
allocation to the analysis populations.

\begin{mermaidfig}
\node[mmin]  (as) at (0,0)       {Assessed for\\eligibility\\(n approx.\ 245\\screened)};
\node[mmdark] (ex) at (7.4,0)    {Excluded: not meeting \S5.1 /\\meeting \S5.2;\\unresectable on staging};
\node[mmstep](rand) at (0,-2.4)  {Randomized 1:1\\(n = 220, central\\permuted block,\\stratified)};
\node[mmgoal](armA) at (-3.2,-5){Arm A (n approx.\ 110)\\daraxonrasib RP2D +\\LLM robotic Whipple};
\node[mmstep](armB) at (3.2,-5) {Arm B (n approx.\ 110)\\mFOLFIRINOX +\\standard Whipple};
\node[mmdark] (di) at (7.4,-5) {Discontinued\\intervention\\(progression,\\withdrawal, \S7)};
\node[mmgoal](an) at (0,-7.6)   {Analyzed: ITT, mITT,\\per-protocol, safety\\(both arms)};
\draw[mmedge] (as) -- node[above,font=\scriptsize]{exclude} (ex);
\draw[mmedge] (as) -- (rand);
\draw[mmedge] (rand) -- (armA);
\draw[mmedge] (rand) -- (armB);
\draw[mmedge] (armA) -- (an);
\draw[mmedge] (armB) -- (an);
\draw[mmedged] (armB) -- (di);
\draw[mmedged] (armA.south) -- ++(0,-3) -| (di.south);\end{mermaidfig}
\vspace{-0.7cm}
\figcaption{Figure 10. Randomized two-arm CONSORT participant flow. Approximately
245 individuals are assessed for eligibility; after exclusions, 220 are randomized
1:1 by central permuted-block randomization stratified by resectability, KRAS
allele, neoadjuvant therapy, and site, to Arm A (daraxonrasib at the RP2D plus the
LLM-directed robotic Whipple, n approximately 110) or Arm B (modified FOLFIRINOX
plus standard\\high-volume Whipple, n approximately 110). Both arms feed the
intention to treat, modified intention to\\treat, per-protocol, and safety analysis
populations, with a protocol-defined discontinuation branch.}

\clearpage

\subsection{Inclusion Criteria}

To be eligible to participate in this study, an individual must meet all of the
following criteria.

\begin{enumerate}[leftmargin=1.6em,itemsep=2pt,topsep=2pt]
\item Signed and dated informed consent obtained before any study-specific
procedure, including the participant's explicit election or declination of the
Physical AI opt-out, which governs use of the LLM-directed advisory layer for those
allocated to the experimental arm (\S8).
\item Age $\geq 18$ years at the time of consent.
\item Histologically or cytologically confirmed PDAC that is resectable or
borderline-resectable on multidisciplinary radiologic review.
\item A documented KRAS G12 mutation (for example G12D, G12V, or G12R) on a
validated tumor genotyping assay.
\item Eastern Cooperative Oncology Group (ECOG) performance status of 0 or 1.
\item Adequate organ and marrow function on screening laboratories, including
hematologic, hepatic, and renal indices sufficient to undergo major
pancreaticobiliary surgery, receive a study regimen.
\item A measured baseline CA 19-9 value (no threshold imposed; value recorded for
serial assessment).
\item Eligible for daraxonrasib under the broad pan-KRAS criteria of RASolute 302
\cite{rasolute302}.
\item Tumor and vascular anatomy compatible with the eight-arm robotic system on
the preoperative imaging review, including port placement and working volume
adequate for operative tasks, for participants who may receive Arm A.
\end{enumerate}

The synthetic reference exemplar PAT-PDAC-0002 illustrates a representative
eligible participant: a 66-year-old with a 3.2~cm head-of-pancreas tumor abutting
the superior mesenteric vein, KRAS G12D, microsatellite stable, CA 19-9 elevated,
ECOG performance status 1, after a completed neoadjuvant modified FOLFIRINOX course
\cite{Conroy2018}. This exemplar is synthetic; it is used for design illustration
and does not represent an enrolled individual.

\subsection{Exclusion Criteria}

An individual who meets any one of the following criteria is excluded from
participation.

\begin{enumerate}[leftmargin=1.6em,itemsep=2pt,topsep=2pt]
\item Distant metastatic disease on staging.
\item Vascular involvement that precludes a margin-negative (R0) resection.
\item Prior pancreatic resection or prior surgery that precludes the planned
reconstruction.
\item Uncontrolled intercurrent illness or comorbidity (ex.: uncontrolled
cardiac, pulmonary, hepatic, renal, or active infectious disease) that would make
major surgery or either study regimen unsafe.
\item Pregnancy or lactation, or unwillingness to use adequate contraception
during study if applicable.
\item Known contraindication or hypersensitivity to daraxonrasib or to modified
FOLFIRINOX or any of their required components, or a concomitant medication
interaction that cannot be safely managed.
\item Inability to provide informed consent or to comply with the Schedule of
Activities.
\end{enumerate}

Participant selection is equitable. Eligibility is determined solely by the
scientific and safety criteria above, and no eligible individual is excluded on the
basis of sex, gender, race, ethnicity, socioeconomic status, or insurance status.
Limited English proficiency is not an exclusion criterion; qualified interpreters
and translated consent materials are provided so that individuals with limited
English proficiency can be fully informed and can participate on the same basis as
English-speaking individuals. The Patient Access and Equity Fund (the
Patient-Aligned Co-Investment Facility, PACIF) removes travel, lodging, lost-wage,
caregiver, and navigation barriers behind a capital firewall, so that
representative enrollment across the eight centers is not limited by a
participant's financial circumstances and so that funders cannot influence
randomization or outcomes \cite{kawchak2026hr9510,ecfr2026part54}.

\subsection{Lifestyle Considerations}

Perioperative lifestyle restrictions apply and are reviewed at consent and at the
baseline visit. Participants follow a standardized perioperative nutrition plan,
including the institutional preoperative fasting protocol and a staged
postoperative diet advancement keyed to recovery of gastrointestinal function and,
in Arm A, to fistula status. Alcohol is to be avoided during the perioperative
period and during any interval of active daraxonrasib dosing, both to limit
hepatotoxic burden and to avoid confounding the safety assessment. Tobacco
cessation is strongly advised and is supported with referral, given its effect on
wound healing and pulmonary risk. No restriction on routine physical activity is
imposed beyond standard postoperative surgical guidance. There are no
study-specific dietary caffeine or grapefruit restrictions beyond those required by
the daraxonrasib labeling and concomitant-medication rules in \S6.

\subsection{Screen Failures}

A screen failure is an individual who consents to participate but is not
subsequently enrolled and randomized because one or more eligibility criteria are
not met. For each screen failure, a minimal CONSORT screen-failure data set is
recorded: demography, the specific criterion or criteria that were not met (the
screen-failure detail), the eligibility determination, and any serious adverse
event that occurred between consent and the determination. No other study
procedures are performed after a screen-failure determination. An individual who
fails screening may be rescreened if the disqualifying condition was transient and
is expected to resolve (for example a correctable laboratory abnormality); a
rescreened individual retains the same participant number, and the rescreening is
documented so that the screening-to-randomization accounting in Figure 10 remains
complete.

\subsection{Strategies for Recruitment and Retention}

Recruitment is conducted across eight academic high-volume HPB centers operating as
a coordinated network. Eligible individuals are identified at each site's tumor-board
review and at surgical and medical-oncology clinics, and the study is described to
them by a member of the research team who is not the treating surgeon, to reduce the
potential for undue influence. The planned accrual rate is approximately 10 to 14
participants per month network-wide over an enrollment period of approximately 30
months, sufficient to randomize 220 participants from approximately 245 screened.

Retention through the 24-month overall-survival window is supported by structured
follow-up with appointment reminders, flexible scheduling, and coordination with
referring providers for participants who travel. Retention is further protected by
the Patient Access and Equity Fund, which provides travel, lodging, lost-wage,
caregiver support, and navigation behind a capital firewall: by removing the
financial and logistical barriers that drive differential dropout, the fund lowers
attrition and thereby protects the statistical power of the comparison and supports
equitable, representative enrollment, while the firewall ensures that funders cannot
influence randomization, endpoint definition, outcome adjudication, or analysis
\cite{kawchak2026hr9510,ecfr2026part54}. 

\clearpage

Because the population is vulnerable,
additional safeguards consistent with Office for Human Research Protections (OHRP)
guidance are applied: enrollment uses an independent consent monitor when indicated,
the consent process explicitly addresses therapeutic misconception about the
LLM-directed robotic system and the voluntariness of the Physical AI opt-out, and
participants are reminded at each contact that they may withdraw at any time without
affecting their clinical care. Compensation is limited to reimbursement of
documented study-related travel, lodging, lost-wage, caregiver, and related expenses
provided through the fund; no payment is offered for undergoing the surgical
procedure or the investigational drug, so that compensation does not function as an
undue inducement.
